\documentclass[11pt,a4paper]{article}
\usepackage{../pelm}
\begin{document}
\answerpaper{Quiz 4}{Weeks 7--8: Grounding, and Verification}

\begin{enumerate}[leftmargin=*, itemsep=0.5em, label=\textbf{\arabic*.}]
\key{B}{Grounding places passages into the context. No retraining is involved --- the model is
  unchanged; only what it can see has changed.}
\key{B}{Retrieve, augment the context, generate. The phrase describes nothing more elaborate than
  that.}
\key{C}{Grounding changes \emph{where} the answer comes from. If the source is wrong or out of date,
  you now have a confident wrong answer with a citation attached --- arguably worse, because the
  citation makes it more persuasive.}
\key{B}{A short, fixed, visually obvious string survives skim-reading. ``I could not find specific
  information, however, based on general principles\ldots'' is the same admission in a form you will
  read straight past.}
\key{C}{No quotation is the signature of answering from memory. If there is no quote, assume memory
  until shown otherwise.}
\key{B}{Over-reading: the claim is stronger than the sentence supports. ``May'' becoming ``must'' is
  the classic instance. Compare claim strength against quote strength.}
\key{B}{Seven answerable --- two easy, three medium, two hard --- plus three plausible unanswerable
  ones.}
\key{B}{In normal use you cannot tell which of your questions are unanswerable, so you never observe
  this behaviour. Constructing it deliberately is the only way to see it.}
\key{B}{Writing the expected answer afterwards allows you to accept whatever you got. The order
  makes the exercise honest.}
\key{C}{Verification requires something outside the system. Repeating a question measures how common
  a pattern is; asking the model about its own confidence stays inside the system.}
\end{enumerate}
\end{document}
